% Auto-generated by generate_assets.py --- tab_position.tex
\begin{table}[htbp]
  \centering
  \small
  \caption{Décomposition du biais de position (en \,\% des paires) : le juge désigne la même réponse dans les deux ordres, choisit toujours la position A, toujours la position B, ou produit un autre motif.}
  \label{tab-position}
  \sisetup{table-number-alignment=center}
  \begin{tabular}{lSSSSS}
    \toprule
    \textbf{Juge} & {$n$} & {\textbf{Cohérent}} & {\textbf{Toujours A}} & {\textbf{Toujours B}} & {\textbf{Autre}} \\
    \midrule
    DeepSeek & 439 & 45.6 & 41.9 & 10.3 & 2.3 \\
    Gemma & 439 & 53.3 & 26.9 & 19.8 & 0.0 \\
    Qwen3 & 439 & 70.8 & 21.6 & 7.3 & 0.2 \\
    \bottomrule
  \end{tabular}
  \caption*{\footnotesize « Toujours A » signifie que le juge désigne la première réponse présentée, indépendamment de son contenu. Le RR de position est égal à la colonne « Cohérent ».}
\end{table}

